% TODO make hw stylesheet
\documentclass[11pt]{scrartcl}
\usepackage[a4paper, total={6in, 8in}]{geometry}
\usepackage[usenames,dvipsnames,svgnames]{xcolor}
\usepackage[shortlabels]{enumitem}
\usepackage[framemethod=TikZ]{mdframed}
\usepackage{amsmath,amssymb,amsthm}
\usepackage[colorlinks]{hyperref}

\hypersetup{
	colorlinks=true,
	linkcolor=blue,
	filecolor=magenta,      
	urlcolor=magenta,
	pdftitle={MAT 319 HW}
}

\usepackage{microtype}
\usepackage{mathtools}
\usepackage[headsepline]{scrlayer-scrpage}
\usepackage{thmtools}
\usepackage{todonotes}

\addtolength{\textheight}{3.14cm}
\providecommand{\ol}{\overline}
\providecommand{\eps}{\varepsilon}
\providecommand{\half}{\frac{1}{2}}
\providecommand{\dang}{\measuredangle} %% Directed angle
\providecommand{\CC}{\mathbb C}
\providecommand{\FF}{\mathbb F}
\providecommand{\NN}{\mathbb N}
\providecommand{\QQ}{\mathbb Q}
\providecommand{\RR}{\mathbb R}
\providecommand{\ZZ}{\mathbb Z}
\providecommand{\dg}{^\circ}
\providecommand{\ii}{\item}
\providecommand{\alert}[1]{{\sffamily\textbf{\textcolor{blue}{#1}}}}
\providecommand{\vect}[1]{\overrightarrow{#1}}
\providecommand{\norm}[1]{\left\|#1\right\|}

\reversemarginpar

\mdfdefinestyle{mdbluebox}{roundcorner=10pt,innerbottommargin=9pt,
	linecolor=blue,backgroundcolor=TealBlue!5,}
\declaretheoremstyle[headfont=\sffamily\bfseries\color{MidnightBlue},
mdframed={style=mdbluebox},]{thmbluebox}

\mdfdefinestyle{mdbloobox}{frametitlefont=\bfseries,innerbottommargin=8pt,
	nobreak=true,backgroundcolor=TealBlue!5,linecolor=blue,}
\declaretheoremstyle[headfont=\bfseries\color{MidnightBlue},
mdframed={style=mdbloobox},headpunct={\\[3pt]},postheadspace=0pt,]{thmbloobox}

\mdfdefinestyle{mdredbox}{frametitlefont=\bfseries,innerbottommargin=8pt,
	nobreak=true,backgroundcolor=Salmon!5,linecolor=RawSienna,}
\declaretheoremstyle[headfont=\bfseries\color{RawSienna},
mdframed={style=mdredbox},headpunct={\\[3pt]},postheadspace=0pt,]{thmredbox}

\mdfdefinestyle{mdgreenbox}{linecolor=ForestGreen,backgroundcolor=ForestGreen!5,
	linewidth=2pt,rightline=false,leftline=true,topline=false,bottomline=false,}
\declaretheoremstyle[headfont=\bfseries\sffamily\color{ForestGreen!70!black},
mdframed={style=mdgreenbox},headpunct={ --- },]{thmgreenbox}

\mdfdefinestyle{mdorangebox}{linecolor=RedOrange,backgroundcolor=RedOrange!5,
	linewidth=2pt,rightline=false,leftline=true,topline=false,bottomline=false,}
\declaretheoremstyle[headfont=\bfseries\sffamily\color{RedOrange!70!black},
mdframed={style=mdorangebox},headpunct={ --- },]{thmorangebox}

\mdfdefinestyle{mdblackbox}{linecolor=black,backgroundcolor=RedViolet!5!gray!5,
	linewidth=3pt,rightline=false,leftline=true,topline=false,bottomline=false,}
\declaretheoremstyle[mdframed={style=mdblackbox}]{thmblackbox}

\declaretheoremstyle[spaceabove=6pt,spacebelow=6pt]{basehead}

\declaretheorem[style=basehead,name=Problem]{problem}
\declaretheorem[style=basehead,name=Problem,numbered=no]{problem*}
\declaretheorem[style=thmbloobox,name=Setup]{setup} % for config geo
\declaretheorem[style=thmredbox,name=Example,sibling=problem]{example}
\declaretheorem[style=thmredbox,name=$\bigstar$ Problem,sibling=problem]{niceprob}
\declaretheorem[style=thmbluebox,name=Theorem,sibling=problem]{theorem}
\declaretheorem[style=thmbluebox,name=Corollary,sibling=theorem]{corollary}
\declaretheorem[style=thmbluebox,name=Proposition,sibling=theorem]{prop}
\declaretheorem[style=thmbluebox,name=Lemma,numberwithin=problem]{lemma}
\declaretheorem[style=thmbluebox,name=Theorem,numbered=no]{theorem*}
\declaretheorem[style=thmbluebox,name=Lemma,numbered=no]{lemma*}
\declaretheorem[style=thmgreenbox,name=Claim,numberwithin=problem]{claim}
\declaretheorem[style=thmgreenbox,name=Claim,numbered=no]{claim*}
\declaretheorem[style=thmblackbox,name=Remark,numberwithin=problem]{remark}
\declaretheorem[style=thmblackbox,name=Remark,numbered=no]{remark*}
\declaretheorem[style=thmgreenbox,name=Definition,sibling=theorem]{definition}
\declaretheorem[style=thmgreenbox,name=Definition,numbered=no]{definition*}
\declaretheorem[style=thmorangebox,name=Warning,sibling=theorem]{warning}
\declaretheorem[style=thmorangebox,name=Warning,numbered=no]{warning*}
\declaretheorem[style=thmorangebox,name=Note,numbered=no]{note}
\declaretheorem[style=thmblackbox,name=Exercise,sibling=problem]{exercise}
\declaretheorem[style=thmblackbox,name=Exercise,numbered=no]{exercise*}
\declaretheorem[style=thmblackbox,name=Question,sibling=problem]{question}
\declaretheorem[style=thmblackbox,name=Question,numbered=no]{question*}

\newenvironment{soln}{\begin{proof}[Solution]}{\end{proof}}
\newenvironment{walkthrough}{\noindent \textbf{Walkthrough.}}{\medskip}
\newlist{walk}{enumerate}{3}
\setlist[walk]{label=\bfseries (\alph*)}

\providecommand{\pow}{\operatorname{Pow}}
\providecommand{\vocab}[1]{{\textbf{\textcolor{ForestGreen}{#1}}}}
\providecommand{\lcm}{\operatorname{lcm}}

\usepackage{asymptote}
\begin{asydef}
	defaultpen(fontsize(10pt));
	unitsize(1cm);
	size(8cm); // set a reasonable default
	usepackage("amsmath");
	usepackage("amssymb");
	settings.tex="pdflatex";
	settings.outformat="pdf";
	// Replacement for olympiad+cse5 which is not standard
	import geometry;
	// recalibrate fill and filldraw for conics
	void filldraw(picture pic = currentpicture, conic g, pen fillpen=defaultpen, pen drawpen=defaultpen)
	{ filldraw(pic, (path) g, fillpen, drawpen); }
	void fill(picture pic = currentpicture, conic g, pen p=defaultpen)
	{ filldraw(pic, (path) g, p); }
	// some geometry
	pair foot(pair P, pair A, pair B) { return foot(triangle(A,B,P).VC); }
	pair orthocenter(pair A, pair B, pair C) { return orthocentercenter(A,B,C); }
	pair centroid(pair A, pair B, pair C) { return (A+B+C)/3; }
	// cse5 abbreviations
	path CP(pair P, pair A) { return circle(P, abs(A-P)); }
	path CR(pair P, real r) { return circle(P, r); }
	pair IP(path p, path q) { return intersectionpoints(p,q)[0]; }
	pair OP(path p, path q) { return intersectionpoints(p,q)[1]; }
	path Line(pair A, pair B, real a=0.6, real b=a) { return (a*(A-B)+A)--(b*(B-A)+B); }
	// cse5 more useful functions
	picture CC() {
		picture p=rotate(0)*currentpicture;
		currentpicture.erase();
		return p;
	}
	pair MP(Label s, pair A, pair B = plain.S, pen p = defaultpen) {
		Label L = s;
		L.s = "$"+s.s+"$";
		label(L, A, B, p);
		return A;
	}
	pair Drawing(Label s = "", pair A, pair B = plain.S, pen p = defaultpen) {
		dot(MP(s, A, B, p), p);
		return A;
	}
	path Drawing(path g, pen p = defaultpen, arrowbar ar = None) {
		draw(g, p, ar);
		return g;
	}
\end{asydef}

\usepackage{mathrsfs}

\ihead{\footnotesize MAT 307 HW01}
\ohead{\footnotesize \today}

\title{\vspace{-2em}MAT 307 HW01}
\author{\large Aditya Pahuja}
\date{\large Due 9/12}
\begin{document}
\maketitle

\begin{problem}
	Using the dot product, prove the Pythagorean theorem.
\end{problem}
\begin{soln}
	Let $\triangle ABC$ be a triangle with side lengths
	$BC = a$, $CA = b$, $AB = c$. We see that
	\[AB^2 = \overrightarrow{AB}\cdot\overrightarrow{AB}
	= (\overrightarrow{AC} + \overrightarrow{CB})\cdot
	(\vect{AC} + \vect{CB}) =
	AC^2 + BC^2 + 2\vect{AC}\cdot\vect{CB}. \]
	That is,
	\[c^2 = a^2 + b^2 + 2\vect{AC}\cdot\vect{CB}.\]
	This immediately shows us that $a^2 + b^2 = c^2$
	if and only if $2\vect{AC}\cdot\vect{CB} = 0$,
	i.e. $AC$ and $BC$ are perpendicular.
\end{soln}

\begin{problem}
	Let $\vec u$ and $\vec v$ be vectors in $\RR^2$
	and recall that the area of the triangle defined by $\vec u$ and $\vec v$
	is given by $A(\vec u, \vec v) = \half|\vec v^\perp \cdot \vec u|$,
	where $\vec v^\perp$ is the counterclockwise rotation of $\vec v$
	by $90\dg$.
	\begin{enumerate}[(a)]
		\ii Show that $(\vec v^\perp\cdot\vec u)^2 =
		\norm{\vec u}^2 \norm{\vec v}^2 - (\vec u\cdot\vec v)^2$.
		Explain what this means geometrically.
		\ii Use this fact to establish Heron's formula.
		\ii Show that, for fixed side lengths $a$ and $b$,
		the triangle with the largest area is right.
		\ii Show that, for a fixed perimeter, the triangle with
		the largest area is equilateral.
	\end{enumerate}
\end{problem}

\begin{soln}
	The proof of (a) is obvious. One way to think about this geometrically
	is to see that the projections of $\vec u$ on $\vec v$
	and $\vec v^\perp$ are the legs of a right triangle with hypotenuse
	$\vec u$; this can be seen by the equivalent equation
	\[\norm{\frac{\vec v\cdot\vec u}{\|\vec v\|^2}\cdot\vec v}^2
	+ \norm{\frac{\vec v^\perp\cdot\vec u}
	{\|\vec v^\perp\|^2}\cdot\vec v^\perp}^2
	= \norm{\operatorname{Proj}_{\vec v} \vec u}^2 +
	\norm{\operatorname{Proj}_{\vec v^\perp} \vec u}^2
	= \|\vec u\|^2\]
	due to the Pythagorean theorem.
	
	(b) is an algebra bash.
	
	(c) follows from (a), since we want to minimize the square of dot product
	of $\vec u$ and $\vec v$, which happens when the dot product is zero
	and thus the vectors are orthogonal.
	
	(d) Jensen carry ngl
\end{soln}

\begin{problem}
	\begin{enumerate}[(a)]
		\ii Let $\vec u,\vec v\in\RR^2$ be non-parallel vectors.
		Describe the set of vectors $s\vec u + t\vec v$ where $s+t = 1$.
		What about if $s$ and $t$ are also nonnegative?
		\ii Let $\vec u, \vec v, \vec w\in\RR^3$ be non-coplanar vectors.
		Describe the set of vectors $r\vec u + s\vec v + t\vec w$
		where $r + s + t = 1$. What happens if $r$, $s$, $t$ are nonnegative?
	\end{enumerate}
\end{problem}

\begin{soln}
	(a): Eliminate $t$ by parameterizing as
	$\vec v + s(\vec u - \vec v)$ to see that this set is the line
	through the endpoints of the vectors;
	if $0\le s\le 1$ then it's just the segment between the endpoints of
	the two vectors.
	
	(b): Same idea, but with ``plane'' instead of ``line''
	and ``triangle'' instead of ``segment.''
\end{soln}

\begin{problem}
	Four vectors are erected perpendicularly to the four faces
	of an arbitrary tetrahedron. Each vector is pointing outward and has length
	equal to the area of the corresponding face.
	Show that these vectors sum to the zero vector.
\end{problem}

\begin{soln}
	Scale the vectors by a factor of two;
	this doesn't affect whether or not they sum to zero,
	so it's equivalent to show that the scaled vectors sum to zero.
	
	Anchor one of the vertices of the tetrahedron to the origin
	and let $\vec a$, $\vec b$, and $\vec c$ be the vectors
	along the edges of the vertices emanating from the origin,
	labeled so that $\vec a\times\vec b$, $\vec b\times\vec c$,
	and $\vec c\times\vec a$ are all pointing outward from the tetrahedron.
	Then, the four vectors are the aforementioned cross products
	along with the product $(\vec c - \vec a) \times (\vec b - \vec a)$.
	The desired sum is then
	\[\vec a \times \vec b + \vec b \times \vec c + \vec c \times \vec a
	+ (\vec c \times \vec b - \vec c \times \vec a - \vec a \times \vec b
	+ \vec a \times \vec a),\]
	using the fact that the cross product distributes over vector addition.
	Then, since the cross product is anticommutative and
	the cross product of a vector with itself is zero (area of parallelogram
	spanned by $\vec a$ and $\vec a$ is zero), we simplify as
	\[\vec b \times \vec c + \vec c \times \vec b + \vec a \times \vec a =
	\vec b\times \vec c - \vec b \times \vec c + \vec 0 = \vec 0.\]
\end{soln}

\begin{problem}
	\begin{enumerate}[(a)]
		\ii Let $\vec u$, $\vec v$, and $\vec w$ be three vectors
		that are not coplanar, i.e.
		$\alpha\vec u + \beta\vec v + \gamma\vec w = 0$
		implies that $\alpha = \beta = \gamma = 0$.
		Show that $\vec u\times\vec v$, $\vec v\times\vec w$, and
		$\vec w\times\vec u$ are not coplanar.
		\ii Suppose $a$, $b$, and $c$ are reals.
		Express the intersection point of planes
		\[\vec u\cdot(x, y, z) = a,\qquad
		\vec v\cdot(x, y, z) = b,\qquad
		\vec w\cdot(x, y, z) = c\]
		as a linear combination of $\vec u\times\vec v$,
		$\vec v\times\vec w$, and $\vec w\times\vec u$.
	\end{enumerate}
\end{problem}

\begin{soln}
	(a) is easy to do by contradiction:
	if the cross products are coplanar, then one of them
	(say $\vec v\times\vec w$) can be written as a linear combination of the
	other two. The other two are both orthogonal to $\vec u$,
	so $\vec v\times\vec w$, being a linear combination of the other two,
	is also orthogonal to $\vec u$; however, this means that
	$\vec u$, $\vec v$, and $\vec w$ are coplanar, contradiction!
	
	(b) Plugging in
	\[(x, y, z) = \alpha \vec v\times\vec w + \beta\vec w\times\vec u
	+ \gamma\vec u\times\vec v,\]
	the three plane equations become
	\begin{align*}
		\vec u\cdot (\alpha \vec v\times\vec w + \beta\vec w\times\vec u
		+ \gamma\vec u\times\vec v) = \vec u\cdot(\alpha\vec v\times \vec w)
		= \alpha(\vec u\cdot(\vec v\times \vec w)) &= a\\
		\vec v\cdot (\alpha \vec v\times\vec w + \beta\vec w\times\vec u
		+ \gamma\vec u\times\vec v) = \vec v\cdot(\beta\vec w\times \vec u)
		= \beta(\vec v\cdot(\vec w\times \vec u)) &= b\\
		\vec w\cdot (\alpha \vec v\times\vec w + \beta\vec w\times\vec u
		+ \gamma\vec u\times\vec v) = \vec w\cdot(\gamma\vec u\times \vec v)
		= \gamma(\vec w\cdot(\vec u\times \vec v)) &= c
	\end{align*}
	after distributing the dot product and noting that
	$\vec x$ is perpendicular to the cross product of $\vec x$
	with any other vector.
	Then,
	\[\vec u\cdot(\vec v\times\vec w) =
	\vec v\cdot(\vec w\times\vec u) =
	\vec w\cdot(\vec u\times\vec v) = D\]
	is the signed area of the parallelepiped formed by $\vec u$, $\vec v$,
	$\vec w$, and we obtain
	$(\alpha, \beta, \gamma) = (a/D, b/D, c/D)$.
\end{soln}

\begin{problem}
	Consider a pyramid in $\RR^n$ with vertices at the origin $O$
	and $\hat e_1$, $\hat e_2$, \dots, $\hat e_n$
	where the $\hat e_i$ form the standard Euclidean basis.
	\begin{enumerate}[(a)]
		\ii Find the coordinates of the point $C$ in the base $B$
		that is equidistant from each of the $\hat e_i$
		and compute $\|\vect{OC}\|$.
		\ii Use Cauchy to show that $C$ is the closest point of $B$ to $O$.
		\ii Compute the angle $\theta$ between $\vect{OC}$ and any $\vect{OV_i}$
		where $V_i$ is the vertex corresponding to $\hat e_i$.
		What happens to $\theta$ as $n$ goes to infinity?
	\end{enumerate}
\end{problem}

\begin{soln}
	(a): Since $C$ is the intersection of
	the perpendicular bisectors of the edges of $B$ (within the plane of $B$),
	it is unique if it exists. I claim that $C = (1/n, 1/n, \dots, 1/n)$
	is the average of $\hat e_1$, $\hat e_2$, \dots, $\hat e_n$;
	this is trivial to check by just computing the length of each $CV_i$.
	
	(b): Let $\vec v$ be a vector with endpoint in $B$ and $\vec c = \vect{OC}$.
	Note that $\|\vec c\| = \sqrt{(1/n)^2\cdot n} = \frac{\sqrt n}{n}$.
	Cauchy-Schwarz shows that
	\[\frac{\|\vec v\|}{\sqrt n} = \|\vec v\| \|\vec c\| \ge
	|\vec v\cdot\vec c| = \left|\frac{x_1}{n} + \cdots + \frac{x_n}{n}\right|
	= \frac 1n,\]
	which means that $\|\vec v\| \ge \frac{\sqrt n}n = \|\vec c\|$
	as desired.
	
	(c): We know that
	\[\frac{\vect{OC}\cdot\vect{OV_i}}{\|\vect{OC}\|\|\vect{OV_i}\|}
	= \cos\theta\]
	where $\theta$ is the angle we are trying to find.
	The left-hand side evaluates to
	\[\frac{(1/n, 1/n, \dots, 1/n)\cdot(0, \dots, 0, 1, 0, \dots, 0)}
	{\sqrt{(1/n)^2\cdot n}\cdot 1}
	= \frac{1/n}{\sqrt n/n} = \frac{\sqrt n}{n},\]
	so \framebox{$\theta = \arccos(\sqrt n/n)$}. Since $\sqrt n/n$
	goes to zero from the positive direction
	as $n$ grows large, $\theta$ approaches $\arccos(0) = \boxed{\pi/2}$.
\end{soln}

\end{document}